\documentclass[a4paper,11pt]{article}

\usepackage[a4paper,margin=2cm]{geometry}
\usepackage[T1]{fontenc}
\usepackage[utf8]{inputenc}
\usepackage[ngerman,english]{babel}
\usepackage{lmodern}
\usepackage{microtype}
\usepackage{xcolor}
\usepackage{parskip}
\usepackage{enumitem}
\usepackage[most]{tcolorbox}
\usepackage{titlesec}
\usepackage[hidelinks]{hyperref}
\usepackage{array}
\usepackage{longtable}
\usepackage[table]{xcolor}

\definecolor{HeaderBlueDark}{RGB}{70,110,170}
\definecolor{RowBlueLight}{RGB}{235,243,255}
\definecolor{RowBlueDark}{RGB}{220,233,250}
\definecolor{SoftGray}{RGB}{90,90,90}

\setlength{\parindent}{0pt}
\setlist[itemize]{leftmargin=1.4em,itemsep=0.35em,topsep=0.35em}
\linespread{1.06}
\renewcommand{\arraystretch}{1.35}
\setlength{\tabcolsep}{6pt}

\titleformat{\section}[block]
{\Large\bfseries\color{HeaderBlueDark}}
{}{0pt}{}
\titlespacing{\section}{0pt}{1.3em}{0.8em}

\titleformat{\subsection}[block]
{\large\bfseries\color{HeaderBlueDark}}
{}{0pt}{}
\titlespacing{\subsection}{0pt}{1.0em}{0.6em}

\newtcolorbox{exbox}{enhanced,breakable,colback=RowBlueLight,colframe=RowBlueDark,boxrule=0.4pt,arc=1.5mm,left=1.2mm,right=1.2mm,top=1mm,bottom=1mm,title={Beispiele \textnormal{(Examples)}},fonttitle=\bfseries,colbacktitle=RowBlueDark,coltitle=black}

\NewDocumentEnvironment{grammarentry}{mm+b}{%
    \begin{tcolorbox}[enhanced,breakable,colback=white,colframe=HeaderBlueDark,boxrule=0.6pt,arc=1.5mm,left=1.5mm,right=1.5mm,top=1mm,bottom=1mm,colbacktitle=HeaderBlueDark,coltitle=white,fonttitle=\bfseries,title={#1\hfill\normalfont\footnotesize #2}]
    #3
    \end{tcolorbox}
}{}

\newcommand{\GermanText}[1]{\textbf{Deutsch: }#1\par}
\newcommand{\EnglishText}[1]{{\small\color{SoftGray}\textbf{English: }#1\par}}
\newcommand{\ExampleItem}[2]{\item \textbf{#1}\\{\small\color{SoftGray}#2}}

\title{Verlauf bis zu einem vergangenen Zeitpunkt\\\large\textnormal{Past Perfect Continuous in German}}
\date{}

\begin{document}
    \maketitle
    \tableofcontents
    \newpage


\section{Grundidee \textnormal{(Basic idea)}}
\begin{grammarentry}{Verlauf statt eigener Tempusform}{No separate German continuous tense}
\GermanText{Das englische Past Perfect Continuous (had been + -ing) hat im Deutschen keine besondere Zeitform. Präteritum + seit kann eine bis damals andauernde Tätigkeit besonders klar beschreiben; Plusquamperfekt kann eine frühere Dauer hervorheben.}
\EnglishText{English had been + -ing has no dedicated German tense form. Preterite + seit clearly describes an activity still continuing at a past reference point; the pluperfect may emphasize the earlier period.}
\end{grammarentry}

\section{Vergangenheit mit seit \textnormal{(Past ongoing actions with seit)}}
\begin{grammarentry}{Damals schon seit ...}{Already ongoing at a past point}
\GermanText{Wenn eine Handlung vor einem Zeitpunkt in der Vergangenheit begonnen hat und damals noch andauerte, ist eine Form mit Präteritum und seit oft besonders eindeutig. Das deutsche Plusquamperfekt ist dafür nicht zwingend erforderlich.}
\EnglishText{When an action began before a past reference point and was still ongoing at that point, a preterite form with seit is often especially clear. German does not necessarily require the pluperfect.}
\end{grammarentry}
\begin{exbox}
\begin{itemize}
\ExampleItem{Als sie anrief, arbeitete ich schon seit zwei Stunden.}{When she called, I had been working for two hours.}
\ExampleItem{Als er endlich kam, warteten wir schon seit einer Stunde.}{When he finally arrived, we had been waiting for an hour.}
\ExampleItem{Damals wohnte sie seit drei Jahren in Graz.}{At that time she had been living in Graz for three years.}
\end{itemize}
\end{exbox}

\section{Past Perfect Continuous \textnormal{(English past perfect continuous)}}
\begin{grammarentry}{Kontext, Präteritum oder Plusquamperfekt}{No mechanical tense equivalence}
\GermanText{Englisches had been + -ing beschreibt häufig eine Dauer bis zu einem vergangenen Zeitpunkt. Deutsch kann dies mit Präteritum + seit oder mit Plusquamperfekt und Zeitangaben ausdrücken. Ein Plusquamperfekt betont häufig die vorhergehende Zeitspanne, beweist aber allein nicht, dass die Handlung zu diesem Zeitpunkt noch lief.}
\EnglishText{English had been + -ing often describes a duration leading up to a past time. German may express it with preterite plus seit, or with the pluperfect and time expressions. The pluperfect often focuses on the preceding period, but by itself does not prove that the activity was still in progress at that moment.}
\end{grammarentry}
\begin{exbox}
\begin{itemize}
\ExampleItem{Als die Sitzung begann, wartete ich schon seit zwanzig Minuten.}{When the meeting started, I had been waiting for twenty minutes.}
\ExampleItem{Als die Sitzung begann, hatte ich schon zwanzig Minuten gewartet.}{By the time the meeting began, I had already waited for twenty minutes; the earlier period is emphasized.}
\ExampleItem{Ich war seit Stunden dabei, die Daten zu überprüfen.}{I had been checking the data for hours / I was in the process of checking them for hours.}
\end{itemize}
\end{exbox}

\section{Zusammenfassung \textnormal{(Summary)}}
\begin{grammarentry}{Merksätze}{Key rules}
\GermanText{Bei einer damals noch laufenden Tätigkeit: Als er kam, wartete ich schon seit einer Stunde. Hatte ich schon eine Stunde gewartet hebt den davorliegenden Zeitraum hervor, sagt aber nicht zwangsläufig dasselbe über das Ende der Tätigkeit.}
\EnglishText{For an activity continuing at the past reference point: Als er kam, wartete ich schon seit einer Stunde. Hatte ich schon eine Stunde gewartet focuses on the earlier period and does not necessarily imply the same endpoint.}
\end{grammarentry}

\section{Quelle \textnormal{(Source)}}
        \GermanText{Institut für Deutsche Sprache (IDS), grammis: Darf man Ich bin am Schreiben schreiben?}
        \EnglishText{Institute for the German Language (IDS), grammis: discussion of progressive constructions and their formation in different tenses.}
        \url{https://grammis.ids-mannheim.de/fragen/4551}
\end{document}
